\section{Rational invariants of a parameter-preserving symplectic lift}
\label{sec:symplectic-lift}

Let $V\in\C[t,a]$ have $\deg_tV=m\ge3$. On affine four-space with
coordinates $(t,s,a,r)$, consider the map $\widehat F$ defined by
\begin{equation}
 a'=a,\qquad s'=s+V_t(t,a),\qquad
 t'=t+s',\qquad r'=r+V_a(t,a).
 \label{eq:symplectic-lift}
\end{equation}
Here $r$ is the fourth coordinate, not the reflection from
Section~\ref{sec:univariate-rigidity}. The preserved parameter $a$ is
already a rational first integral. Our question is whether there are any
others, in the full field $\C(t,s,a,r)$.

The map is a polynomial symplectic automorphism for
\[
 \omega=\dd t\wedge\dd s+\dd a\wedge\dd r.
\]
Indeed, first apply the kick
\[
 (t,s,a,r)\longmapsto
 \bigl(t,s+V_t(t,a),a,r+V_a(t,a)\bigr).
\]
Its pullback changes $\omega$ by
\[
 V_{ta}\,\dd t\wedge\dd a+
 V_{at}\,\dd a\wedge\dd t=0.
\]
The following drift $(t,s,a,r)\mapsto(t+s,s,a,r)$ also preserves $\omega$,
because $\dd(t+s)\wedge\dd s=\dd t\wedge\dd s$. Both steps have polynomial
inverses: subtract the same gradient in the kick, and subtract $s$ from
$t$ in the drift. Their composition is \eqref{eq:symplectic-lift}.

Regard $a$ as a member of the constant field and put
\begin{equation}
 K=\C(a),\qquad x=t,\qquad y=t-s,\qquad
 p(x)=2x+V_t(x,a).
 \label{eq:lift-henon-coordinates}
\end{equation}
Then $s=x-y$, and the induced map on $(x,y)$ is
\[
 x'=p(x)-y,\qquad y'=x.
\]
It is a single H\'enon map $H_p$ over $K$, with
$\deg_xp=m-1\ge2$. Write $L=K(x,y)$ and $\sigma=H_p^*$ for its pullback.
Corollary~\ref{cor:rational-reduction} gives $L^\sigma=K$. On the remaining
transcendental coordinate, the pullback $\widehat\sigma=\widehat F^*$ acts
by
\[
 \widehat\sigma r=r+b,\qquad b=V_a(x,a)\in L,
 \qquad \widehat\sigma|_L=\sigma.
\]

The reduction from an extra invariant to an additive coboundary is the
classical additive-extension criterion of Karr, as recalled by
Schneider~\cite{Schneider2016DifferenceRings}. We give the coefficient
argument, including the full constant field in the solvable case.

\begin{lemma}[Additive-extension criterion]
\label{lem:additive-extension}
Let $L$ be a field of characteristic zero with automorphism $\sigma$, and
let $K=L^\sigma$. For $b\in L$, extend $\sigma$ to an automorphism
$\widehat\sigma$ of $L(r)$, where $r$ is transcendental over $L$, by
$\widehat\sigma r=r+b$. If there is no $h\in L$ with
$\sigma h-h=b$, then
\[
 L(r)^{\widehat\sigma}=K.
\]
If such an $h$ exists, then
\begin{equation}
 L(r)^{\widehat\sigma}=K(r-h).
 \label{eq:additive-fixed-field}
\end{equation}
\end{lemma}

\begin{proof}
Let $I=P/Q$ be a nonzero invariant, where $P,Q\in L[r]$ are nonzero and
coprime. Its invariance says
\[
 (\widehat\sigma P)Q=P(\widehat\sigma Q).
\]
Since $P$ and $Q$ are coprime, $P$ divides $\widehat\sigma P$ and $Q$
divides $\widehat\sigma Q$. The substitution $r\mapsto r+b$ preserves
$r$-degree, so the quotients are nonzero members of $L$. The displayed
identity makes them equal: for some $\lambda\in L^*$,
\begin{equation}
 \widehat\sigma P=\lambda P,\qquad
 \widehat\sigma Q=\lambda Q.
 \label{eq:additive-semiinvariants}
\end{equation}
Let $p_j$ and $q_\ell$ be the leading coefficients of $P$ and $Q$,
respectively. Comparing leading coefficients in
\eqref{eq:additive-semiinvariants} gives
\[
 \sigma p_j=\lambda p_j,\qquad
 \sigma q_\ell=\lambda q_\ell.
\]
Thus the monic polynomials $P/p_j$ and $Q/q_\ell$ are individually fixed by
$\widehat\sigma$. Moreover, the leading-coefficient ratio is fixed in
the base field:
\begin{equation}
 \sigma(p_j/q_\ell)=p_j/q_\ell,
 \qquad p_j/q_\ell\in K.
 \label{eq:leading-coefficient-ratio}
\end{equation}

If $I$ depends on $r$, at least one of the two monic polynomials has
positive degree. Denote it by
\[
 T=r^q+u_{q-1}r^{q-1}+\cdots+u_0,\qquad q\ge1.
\]
The coefficient of $r^{q-1}$ in $\widehat\sigma T=T$ is
\[
 qb+\sigma u_{q-1}=u_{q-1}.
\]
Division by the nonzero integer $q$ yields
$b=\sigma h-h$ with $h=-u_{q-1}/q$. Hence, when no such $h$ exists, every
invariant lies in $L$ and therefore in $L^\sigma=K$.

Conversely, suppose $b=\sigma h-h$ and put $z=r-h$. Then
\[
 \widehat\sigma z=r+b-\sigma h=r-h=z.
\]
The element $z$ is transcendental over $L$, so $L(r)=L(z)$, and $K(z)$ is
contained in the fixed field. To prove the reverse inclusion, write an
arbitrary nonzero invariant as a coprime quotient in $L[z]$. The preceding
divisibility and leading-coefficient argument applies with $z$ in place of
$r$. It gives invariant monic numerator and denominator, and their
leading-coefficient ratio belongs to $K$ by
\eqref{eq:leading-coefficient-ratio}. Because $z$ is fixed, equality of
each monic polynomial with its transform means that every coefficient is
fixed by $\sigma$, hence belongs to $K$. The quotient, including its
leading-coefficient ratio, is therefore in $K(z)$. This proves
\eqref{eq:additive-fixed-field}.
\end{proof}

The next-to-leading-coefficient reduction also occurs in the contact lifts
studied by Cerveau and D\'eserti~\cite{CerveauDeserti2018Contact}. Their
contact setting is distinct from the four-dimensional symplectic map
\eqref{eq:symplectic-lift}; here the specific remaining equation is the
univariate H\'enon coboundary equation classified in
Theorem~\ref{thm:univariate}.

\begin{theorem}[Complete rational fixed field of the lift]
\label{thm:lift-field}
For the potential and map in \eqref{eq:symplectic-lift}, the field of all
rational first integrals is
\begin{equation}
 \C(t,s,a,r)^{\widehat F}=
 \begin{cases}
  \C\bigl(a,r-c(a)s\bigr),
    &\text{if }V_a=c(a)V_t\text{ for some }c\in\C(a),\\
  \C(a),&\text{if no such }c\text{ exists}.
 \end{cases}
 \label{eq:lift-fixed-field}
\end{equation}
\end{theorem}

\begin{proof}
The change \eqref{eq:lift-henon-coordinates} identifies
$\C(t,s,a,r)$ with $L(r)$, and the pullback is the additive extension
described above. Lemma~\ref{lem:additive-extension} reduces the question
to the existence of $h\in K(x,y)$ such that
\[
 \sigma h-h=V_a(x,a).
\]
Theorem~\ref{thm:univariate}, applied over $K=\C(a)$, says precisely that
such an $h$ exists if and only if
\[
 V_a(x,a)=c(a)\bigl(p(x)-2x\bigr)=c(a)V_t(x,a)
\]
for some $c\in K$. When it exists, every primitive is
$h=c(a)(x-y)+c_0(a)=c(a)s+c_0(a)$. The constant $c_0(a)$ does not affect
the field $K(r-h)$, which is $\C(a,r-c(a)s)$. The two alternatives of the
lemma now give \eqref{eq:lift-fixed-field}.
\end{proof}

In particular, \eqref{eq:lift-fixed-field} is not restricted to polynomial
integrals or to a bounded degree. The calculation is over $\C(a)$ and
does not assert that every specialized parameter fiber has an identical
degree or invariant classification.

\begin{proposition}[Polynomial translation exception]
\label{prop:translation-exception}
The condition $V_a=c(a)V_t$ for some $c\in\C(a)$ is equivalent to
\begin{equation}
 V(t,a)=W(t+A(a)),\qquad
 A\in\C[a],\qquad W\in\C[t].
 \label{eq:translation-potential}
\end{equation}
In this case $c=A'$, and the coordinates
\begin{equation}
 u=t+A(a),\qquad \rho=r-A'(a)s,\qquad s,\qquad a
 \label{eq:translation-coordinates}
\end{equation}
give a global polynomial symplectic change of variables. They transform
$\widehat F$ into
\begin{equation}
 s'=s+W'(u),\qquad u'=u+s',\qquad a'=a,\qquad\rho'=\rho.
 \label{eq:decoupled-lift}
\end{equation}
The symbol $\rho$ in this section denotes the coordinate in
\eqref{eq:translation-coordinates}, independently of the finite digit
permutation used in Section~\ref{sec:hilbert-series}.
\end{proposition}

\begin{proof}
Write
\[
 V(t,a)=\sum_{j=0}^m v_j(a)t^j,\qquad
 v_j\in\C[a],\qquad v_m\ne0.
\]
Suppose $V_a=c(a)V_t$. The coefficient of $t^m$ on the right is zero,
whereas on the left it is $v_m'$. Hence $v_m'=0$ and
$v_m\in\C^*$. Comparing the coefficients of $t^{m-1}$ next gives
\[
 v_{m-1}'=m c(a)v_m.
\]
Define
\begin{equation}
 A(a)=\frac{v_{m-1}(a)}{m v_m}\in\C[a].
 \label{eq:translation-from-top-coefficients}
\end{equation}
Then $c=A'$. In particular, the polynomial nature of this antiderivative
follows from the two highest coefficients; it is not an assumption on
the initially rational function $c$.

For a new variable $u$, the polynomial
$\widetilde V(u,a)=V(u-A(a),a)$ satisfies, with $u$ held fixed,
\[
 \frac{\partial\widetilde V}{\partial a}
 =V_a(u-A(a),a)-A'(a)V_t(u-A(a),a)=0.
\]
Since $\widetilde V\in\C[u,a]$ and the characteristic is zero, every
coefficient of a positive power of $a$ vanishes. Thus
$\widetilde V=W(u)$ for a polynomial $W\in\C[u]$, proving
\eqref{eq:translation-potential}. Conversely, differentiating
\eqref{eq:translation-potential} gives
$V_a=A'(a)W'(t+A(a))=A'(a)V_t$.

The inverse of \eqref{eq:translation-coordinates} is polynomial:
\[
 t=u-A(a),\qquad r=\rho+A'(a)s,
\]
with $a$ and $s$ unchanged. Its symplectic property follows from
\[
 \dd u=\dd t+A'(a)\dd a,\qquad
 \dd\rho=\dd r-A'(a)\dd s-A''(a)s\,\dd a.
\]
Indeed,
\begin{align*}
 \dd u\wedge\dd s+\dd a\wedge\dd\rho
 &=\dd t\wedge\dd s+A'(a)\dd a\wedge\dd s
   +\dd a\wedge\dd r-A'(a)\dd a\wedge\dd s\\
 &=\dd t\wedge\dd s+\dd a\wedge\dd r.
\end{align*}
The omitted $A''$ term has the factor $\dd a\wedge\dd a=0$.
Finally, $V_t=W'(u)$ and $V_a=A'W'(u)$. Since $a'=a$, direct substitution
in \eqref{eq:symplectic-lift} yields
\[
 u'=t'+A(a)=u+s',\qquad
 \rho'=r+A'W'(u)-A'\bigl(s+W'(u)\bigr)=\rho,
\]
as well as the other two equations in \eqref{eq:decoupled-lift}.
\end{proof}

Thus the exceptional fixed field can be written as $\C(a,\rho)$, and
the map is a product in the global symplectic coordinates $(u,s,a,\rho)$.
This description requires the exact potential
\eqref{eq:translation-potential}, without an arbitrary added polynomial in
$a$. A nonconstant such polynomial has no effect on the two-dimensional
base map but does change the fourth equation of the lift.

\begin{corollary}[Absence of a rational Liouville pair]
\label{cor:no-liouville-pair}
The map \eqref{eq:symplectic-lift} has no pair of algebraically independent
rational first integrals that Poisson commute for $\omega$.
\end{corollary}

\begin{proof}
Use the convention $\{t,s\}=\{a,r\}=1$, with brackets between the two
coordinate pairs equal to zero. If the second case of
\eqref{eq:lift-fixed-field} holds, the invariant field has transcendence
degree one over $\C$, and so cannot contain two algebraically independent
elements.

In the exceptional case the entire invariant field is $\C(a,\rho)$, and
the symplectic coordinate change gives $\{a,\rho\}=1$. For any two of its
elements $I,J$, the bracket is
\begin{equation}
 \{I,J\}=I_aJ_\rho-I_\rho J_a.
 \label{eq:invariant-field-poisson}
\end{equation}
If $I,J$ are algebraically independent, $\C(a,\rho)$ is a finite
algebraic extension of $\C(I,J)$, and it is separable in characteristic
zero. Therefore $\dd I$ and $\dd J$ remain linearly independent in the
differential space over $\C(a,\rho)$. Equivalently,
\[
 \dd I\wedge\dd J
 =(I_aJ_\rho-I_\rho J_a)\,\dd a\wedge\dd\rho\ne0.
\]
Equation~\eqref{eq:invariant-field-poisson} then gives $\{I,J\}\ne0$.
This excludes every algebraically independent pair in the invariant
field, not only its displayed generators.
\end{proof}

For example, $V=t^3/3+at$ has $V_a=t$ and $V_t=t^2+a$, which are not
proportional over $\C(a)$; its fixed field is $\C(a)$. The translation
potential $V=(t+a^2)^3/3$ has $c=2a$ and the additional invariant
$\rho=r-2as$. In contrast, $V=(t+a^2)^3/3+a$ has the same base map
as the second example, but in the same coordinates $\rho'=\rho+1$.
Indeed, its derivatives satisfy
$V_a=2aV_t+1$, so comparison of the highest $t$ coefficient would force
$c=2a$ and leave the nonzero constant $1$. Its fixed field is therefore
again $\C(a)$. The distinction concerns the whole lift, not a
classification of all base families whose parameter can be removed by
some change of coordinates. The corollary likewise concerns rational
integrals for the stated symplectic structure, not other function classes
or local notions of integrability.

The ordinary-degree filtration connects the global coboundary equation
to exact obstruction counts and to detection on a finite periodic scheme.
For the lift, the same equation determines the entire rational invariant
field and isolates a global polynomial translation exception. The
periodic tests throughout the paper remain identities in complete
fixed-point coordinate algebras. Over an algebraically closed field,
geometric point values give the same test when the relevant algebra is
reduced; without reducedness they only imply that the orbit sum is
nilpotent. Whether the stated long-period hypotheses suffice for a
geometric-point-only test without a reducedness assumption is not
established here.
